% titel: Blatt 5, A6: Vergleich des Ionenradius
% skript: Mortimer Kap. 7 (chemie: Ionenbindung, Ionenradius)

Ionen haben andere Radien als ihre Atome: Kationen sind kleiner (weniger
Elektronen, gleiche Kernladung), Anionen grösser (mehr Elektronen). Bei
isoelektronischen Ionen (gleiche Elektronenzahl) entscheidet die Kernladung.

\begin{werkzeugkasten}[Ionenradius vergleichen]
\begin{itemize}
  \item Kation  $<$ Neutralatom; Anion $>$ Neutralatom (bei gleichem Element).
  \item \textbf{Isoelektronische Reihe} (gleiche Elektronenzahl, z.\ B. alle
        mit $[\mathrm{Ar}]$-Rumpf): je \textit{kleiner} $Z$, desto \textit{grösser}
        der Radius $\to$ stärker negatives Ion ist grösser.
  \item Gleiche Ladung, gleiche Gruppe: je tiefer (mehr Schalen), desto grösser.
  \item Hohe Ladung zieht die Restschale stärker zusammen $\to$ kleineres Ion.
\end{itemize}
\textit{Skript:} Mortimer Kap. 7 (Ionenradius).
\end{werkzeugkasten}

\begin{achtung}[Darauf achten]
\begin{itemize}
  \item Bei Paaren wie $\mathrm{Ca^{2+}}/\mathrm{P^{3-}}$ beide Elektronenzahlen
        zählen: $\mathrm{Ca^{2+}}$ und $\mathrm{P^{3-}}$ sind isoelektronisch
        ($18\,\mathrm{e^-}$) --- dann gewinnt die kleinere Kernladung.
  \item Kationen verschiedener Ladung (z.\ B. $\mathrm{K^+}$ vs.\ $\mathrm{Sc^{3+}}$):
        gleiche Elektronenzahl, kleinere Kernladung $\to$ grösserer Radius.
  \item Nicht die Ladungszahl verwechseln: höhere positive Ladung heisst
        \textit{kleiner}, nicht grösser.
\end{itemize}
\end{achtung}
